% Standalone research fragment: exact coefficient proof and arithmetic witness.
% The two global inputs are the normalized signed-congruence theorem and
% simultaneous rational Kummer surjectivity.  No coefficient identity below
% depends on the accompanying computer calculations.

\section{A cubic Magnus normalization beyond rational radicals}
\label{sec:cubic-magnus-normalization}

Put $G=\widehat{\mathrm{GT}}$, let $H\subseteq G$ be the Galois image, and
write $N=\ker(\lambda,\rho_2^I)$.  For a multiplier-one element
$g=(1,f)$, write its pro-$2$ Magnus expansion as
\[
 x=1+S,\qquad y=1+T,\qquad
 m_3(g)=[S^2T]f\in\mathbf Z_2.
\]
Here brackets around a word mean its coefficient, whereas group commutators
are $[a,b]=aba^{-1}b^{-1}$.  The notation $m_3$ separates this cubic
monodromy coordinate from the rational Kummer coordinate $\rho_3^I$.

We use two previously established global inputs.  First, for $g\in N$,
\begin{equation}\label{eq:magnus-normalized-congruence-input}
 f(U^2,V^2)=I,
 \qquad
 U=\begin{pmatrix}1&1\\0&1\end{pmatrix},\quad
 V=\begin{pmatrix}1&0\\-1&1\end{pmatrix}
\end{equation}
in $\mathrm{SL}_2(\widehat{\mathbf Z})$.  Second, over
$\mathbf Q^{\mathrm{ab}}$, the simultaneous rational Kummer map has image
$\prod_q2\widehat{\mathbf Z}$, with $q$ running over all rational primes.
Thus a fiber of a fixed cyclotomic multiplier is the full coset determined
by the quadratic Kummer characters.  These are the global inputs to the argument below, proved in
Proposition~\ref{prop:normalized-congruence-triviality} and
Theorem~\ref{thm:main-rational-normalization}.

\begin{lemma}[An integral alternating coefficient identity]
\label{lem:integral-alternating-cubic}
Let $f$ be a pro-$2$ commutator word satisfying
\[
 f(y,x)f(x,y)=1,\qquad
 f(z,x)f(y,z)f(x,y)=1,\qquad z=(xy)^{-1}.
\]
Set $a=[S^2T]f$ and $A=[S^2T^2S]f$.  Then
\begin{equation}\label{eq:alternating-cubic-identity}
 [STS]f=-2a,\qquad [ST^2S]f=a,\qquad
 [STSTS]f=-2a-4A.
\end{equation}
In particular, $[STSTS]f\equiv2a\pmod4$.
\end{lemma}

\begin{proof}
The degree-two part of $f$ is a multiple of $[S,T]$.  The degree-two
part of the three-cycle identity is three times that multiple, so it is
zero.  Thus $f\in\gamma_3$ in the pro-$2$ nilpotent quotient relevant
to the calculation.

For completeness, the low-weight collection calculation needed here is
given explicitly.  Let $L_w$ be the standard Lyndon group commutator
associated to a Lyndon word $w$ on $x<y$: $L_x=x$, $L_y=y$, and
$L_w=[L_u,L_v]$ for the standard factorization $w=uv$ with $v$ the
longest proper Lyndon suffix.  Up to weight five, collect $f$ as the
ordered product of these commutators, with coordinates
\[
\begin{gathered}
\begin{array}{c|rrrrr}
 w&xxy&xyy&xxxy&xxyy&xyyy\\
 \hline
 \text{coordinate}&a&d&u&v&w
\end{array}
\\[5pt]
\begin{array}{c|rrrrrr}
 w&xxxxy&xxxyy&xxyxy&xxyyy&xyxyy&xyyyy\\
 \hline
 \text{coordinate}&b&h&i&j&e&k.
\end{array}
\end{gathered}
\]
All coordinates belong to $\mathbf Z_2$.  Since every factor begins in
degree at least three, the product is linear in these coordinates
through degree five.  Write $I=f(y,x)f(x,y)-1$ and
$II=f(z,x)f(y,z)f(x,y)-1$.  Direct substitution of
$z=(1+T)^{-1}(1+S)^{-1}$ gives the following five coefficients:
\[
\begin{array}{c|l}
\text{coefficient}&\text{value}\\ \hline
{[S^2T]I}&a+d\\
{[S^3T]I}&u-w\\
{[S^3T^2]I}&h+j-u+w\\
{[S^3T]II}&a-2d+2u+v+2w\\
{[S^4T]II}&-a+3d-h-j-2u-v-3w.
\end{array}
\]
Each entry is zero.  The first three give $d=-a$, $w=u$, and
$h+j=0$; the last two then give
\[
 3a+4u+v=0,\qquad -4a-5u-v=0.
\]
Consequently $u=w=-a$ and $v=a$.  No division by two has occurred.
The same commutator expansion gives
\begin{equation}\label{eq:cubic-three-small-coefficients}
 [STS]f=-2a,\qquad
 [ST^2S]f=a,\qquad
 [STS^2]f=4a+3u=a.
\end{equation}
This supplies all the low-weight collection identities used below.

Let $C=[ST^2S^2]f$.  The group-ring antipode sends each generator to
its inverse and reverses products.  Since $f-1$ starts in degree three,
$f^{-1}=1-(f-1)$ through degree five.  The coefficient of $S^2T^2S$
in the antipode of $f$ is therefore
\[
 -A=-\bigl(C+[STS^2]f+[ST^2S]f+[STS]f\bigr)=-C.
\]
Thus $C=A$.

Finally, group-likeness of $f$ for
$\Delta(S)=S\otimes1+1\otimes S+S\otimes S$ and the analogous
formula for $T$ gives, on taking the coefficient of
$S\otimes ST^2S$,
\[
 0=2[S^2T^2S]f+[STSTS]f+2[ST^2S^2]f+2[ST^2S]f.
\]
The right-hand side counts the five insertions of the extra $S$ and
the two possible mergers with an existing $S$.  Substitution of
$C=A$ and \eqref{eq:cubic-three-small-coefficients} proves
\eqref{eq:alternating-cubic-identity}.
\end{proof}

\begin{proposition}[Divisibility and the cubic character]
\label{prop:cubic-magnus-eight}
One has
\[
 m_3(N)\subseteq8\mathbf Z_2.
\]
Consequently
\[
 \eta_3:N\longrightarrow\mathbf F_2,
 \qquad \eta_3(g)=m_3(g)/8\pmod2
\]
is an ordinary continuous $G$-invariant quadratic character.
\end{proposition}

\begin{proof}
Under the substitution in
\eqref{eq:magnus-normalized-congruence-input},
\[
 S=\begin{pmatrix}0&2\\0&0\end{pmatrix},\qquad
 T=\begin{pmatrix}0&0\\-2&0\end{pmatrix}.
\]
Consecutive equal letters vanish.  Every even-length surviving word
is diagonal, and every word of degree at least seven vanishes modulo
$128$.  The only contributions to the upper-right entry modulo $128$
are therefore
\[
 [STS]f\,(-8)+[STSTS]f\,32
 =16a+32(-2a-4A)\equiv80a\pmod{128}.
\]
The matrix identity makes this zero.  Since $80=16\cdot5$, it follows
that $a\in8\mathbf Z_2$.

On $\ker\lambda$, the cubic Magnus coordinate is additive under the
GT product.  Moreover, for $s\in G$ and $n\in\ker\lambda$,
\begin{equation}\label{eq:cubic-scalar-conjugation}
 m_3(sns^{-1})=\lambda_2(s)^3m_3(n).
\end{equation}
Indeed, a multiplier-one word has no part below degree three; the
GT product adds its degree-three terms, and conjugation acts on that
degree by the cube of the degree-one scalar.  Terms from the
conjugating word start in degree two and cannot change this leading
degree-three term.  These statements may equally be read from the
induced tangential automorphism in degree four.  Continuity follows
from the Magnus coefficient map.  Finally, every $2$-adic unit has
odd cube, so division by eight and reduction modulo two turn
\eqref{eq:cubic-scalar-conjugation} into invariance.
\end{proof}

We next give the arithmetic witness, including the small congruence
which detects its parity.

\begin{lemma}[A boundary parity congruence]
\label{lem:rho2-cubic-modfour}
For every multiplier-one gentle element,
\[
 \rho_2^I(g)\equiv-2m_3(g)\pmod4
\]
in the $2$-adic coordinate.
\end{lemma}

\begin{proof}
Use the affine matrix representation
\[
 x\longmapsto D=\begin{pmatrix}-1&0\\0&1\end{pmatrix},\qquad
 y\longmapsto E=\begin{pmatrix}1&1\\0&1\end{pmatrix}.
\]
In the index-two cyclic subgroup, $x^2$ maps to the identity, while
$Y_0=y$ and $Y_1=xyx^{-1}$ have translation coordinates $1$ and $-1$.
Since $f$ has zero total $y$-exponent, its translation is
$2\ell_2(f)=-2\rho_2^I(f)$.

On the Magnus side, $S=D-I=\operatorname{diag}(-2,0)$ and $T=E-I$.
Here $T^2=TS=0$.  Since the first nonzero term of $f$ has degree
three, its translation is
\[
 \sum_{n\ge2}(-2)^n[S^nT]f\equiv4m_3(f)\pmod8.
\]
The series converges $2$-adically.  Comparing translations and
dividing the congruence by $-2$ proves the claim.
\end{proof}

Let
\[
 D=N_{\mathrm{rad}}=\bigcap_q\ker(\rho_q^I)\cap\ker\lambda,
 \qquad
 K_{\mathrm{rad}}=D\cap H=G_L,
\]
where $L=\mathbf Q(\mu_\infty,q^{1/n}:q\text{ prime},n\ge1)$.
The all-prime normalization gives $G=DH$.

\begin{proposition}[An arithmetic value outside the rational Kummer sector]
\label{prop:cubic-arithmetic-witness}
There exists $k\in K_{\mathrm{rad}}$ such that
\[
 m_3(k)\in8\mathbf Z_2^\times.
\]
In particular, $\eta_3|_{K_{\mathrm{rad}}}\ne0$, so $\eta_3$ is
independent of every character generated by rational Kummer coordinates.
Moreover,
\[
 m_3(D)=m_3(K_{\mathrm{rad}})=8\mathbf Z_2.
\]
\end{proposition}

\begin{proof}
Simultaneous rational Kummer surjectivity supplies
$\sigma\in H^{\mathrm{cyc}}=G_{\mathbf Q^{\mathrm{ab}}}$ with
\[
 \rho_2^I(\sigma)=2\in\widehat{\mathbf Z},\qquad
 \rho_q^I(\sigma)=0\quad(q\text{ odd}).
\]
Lemma~\ref{lem:rho2-cubic-modfour} shows that $m_3(\sigma)$ is odd.
Choose $h\in H$ with $\lambda_2(h)=3$, and put
$t=\lambda(h)\in\widehat{\mathbf Z}^{\times}$.  For
$a\in\widehat{\mathbf Z}$, the power $\sigma^a$ means the value of
the continuous map $\widehat{\mathbf Z}\to
\overline{\langle\sigma\rangle}$ extending integer powers.
Define
\[
 k=(h\sigma h^{-1})\sigma^{-t}.
\]
Every rational Kummer coordinate of this element is zero:
on the multiplier-one subgroup the cocycles are homomorphisms,
and conjugation by $h$ multiplies each by $t$, exactly as the
profinite power $\sigma^t$ does.  Thus $k\in K_{\mathrm{rad}}$.
By additivity and \eqref{eq:cubic-scalar-conjugation},
\[
 m_3(k)=(3^3-3)m_3(\sigma)
        =24m_3(\sigma)\in8\mathbf Z_2^\times.
\]
Every rational Kummer character vanishes on $K_{\mathrm{rad}}$, while
$\eta_3(k)=1$, proving the independence assertion.

The image of the compact group $K_{\mathrm{rad}}$ under $m_3$ is a
closed additive subgroup of $\mathbf Z_2$.  It is contained in
$8\mathbf Z_2$ by Proposition~\ref{prop:cubic-magnus-eight}, and it
contains a generator of that closed subgroup.  It is therefore all of
$8\mathbf Z_2$.  The same conclusion follows for $D$ by inclusion.
\end{proof}

\begin{corollary}[A further proper normalization]
\label{cor:cubic-radical-normalization}
Define
\[
 D_3=N_{\mathrm{rad},3}=D\cap\ker m_3,
 \qquad K_{\mathrm{rad},3}=K_{\mathrm{rad}}\cap\ker m_3.
\]
Then $D_3$ is a proper closed normal subgroup of $G$, and
\[
 D=D_3K_{\mathrm{rad}},\qquad
 G=D_3H,\qquad D_3\cap H=K_{\mathrm{rad},3}.
\]
Writing
\[
 B=N/\overline{[G,N]K},\qquad K=N\cap H,
\]
the previously obtained radical normalization refines canonically to
\[
 B\simeq
 D/\overline{[G,D]K_{\mathrm{rad}}}
 \simeq
 D_3/\overline{[G,D_3]K_{\mathrm{rad},3}}.
\]
If $L_3=\overline{\mathbf Q}^{\,K_{\mathrm{rad},3}}$, then
$\operatorname{Gal}(L_3/L)\simeq\mathbf Z_2$.  This defines the
arithmetic extension by the explicit cubic monodromy coordinate; no
radical formula for $L_3$ is asserted.
\end{corollary}

\begin{proof}
Closedness and normality follow from continuity and
\eqref{eq:cubic-scalar-conjugation}.  Properness and the factorization
$D=D_3K_{\mathrm{rad}}$ follow from the equality of the two cubic
images.  The remaining factorization and intersection statements are
immediate from $G=DH$.

For the quotient assertion, one may use the following elementary
group argument.  If a profinite group factors as $G=EH$ with $E$
closed normal, then
\[
 E/\overline{[G,E](E\cap H)}
 \simeq G/\overline{[G,E]H}.
\]
The right-hand quotient is abelian because the image of $E$ is
central and it generates after the image of $H$ is killed.  If also
$E=E'(E\cap H)$ and $G=E'H$, then
$\overline{[G,E]H}=\overline{[G,E']H}$: after quotienting by the latter,
the image of $E'$ is central and $H$ is trivial.  Apply this with
$E=D$, $E'=D_3$, and combine with the earlier radical normalization.
Finally, the kernel of the surjective map
$m_3:G_L\to8\mathbf Z_2$ is exactly $G_{L_3}$.
\end{proof}

\begin{corollary}\label{cor:cubic-radical-joint-image}
The map
\[
 N\longrightarrow
 \left(\prod_{q\text{ odd prime}}2\widehat{\mathbf Z}\right)
       \times8\mathbf Z_2,
 \qquad n\longmapsto((\rho_q^I(n))_q,m_3(n))
\]
has kernel $D_3$ and is surjective; already its restriction to
$K=N\cap H$ is surjective.
Consequently the $G$-invariant continuous $\mathbf F_2$-characters
of $N$ which vanish on $D_3$ form exactly
\[
 \left(\bigoplus_{q\text{ odd prime}}\mathbf F_2\psi_q\right)
                \oplus\mathbf F_2\eta_3.
\]
For an odd prime $p$, no nonzero $G$-invariant
$\mathbf F_p$-character of $N$ vanishes on $D_3$.
\end{corollary}

\begin{proof}
Prescribe any rational Kummer vector and cubic value in the displayed
target.  First choose $h\in K$ with the prescribed Kummer vector.
Its cubic coordinate lies in $8\mathbf Z_2$.  The difference between
that coordinate and the prescribed one can be corrected by an element
of $K_{\mathrm{rad}}$, by
Proposition~\ref{prop:cubic-arithmetic-witness}; this changes none of
the Kummer coordinates.  The kernel is exactly the common kernel $D_3$.
Compactness identifies $N/D_3$ with the displayed product.

Every continuous character of this product into a finite group depends
on finitely many coordinates.  At two, each divided coordinate has a
one-dimensional character space, and conjugation multiplies the Kummer
coordinates by $\lambda$ and the cubic coordinate by $\lambda_2^3$.
Both scalars are odd, so every displayed quadratic character is invariant.
At an odd prime $p$, the pro-$2$ cubic factor has no nonzero character
into $\mathbf F_p$.  A Kummer character is multiplied by $\lambda_p$
under conjugation; choosing a multiplier with $\lambda_p\not\equiv1$
modulo $p$ forces every invariant coefficient to vanish.
\end{proof}

The character $\eta_3$ has nonzero arithmetic restriction.  The result
therefore enlarges the explicitly normalized arithmetic character
sector; it does not construct a nonzero character in the kernel of
restriction to the Galois subgroup.
